\section{Background}\label{sec:backgound}

% 两种，一种处理ingest任务，一种处理retreivial任务。前者要低时延，后者要高吞吐。
% 数据video或者image，都是rgb space。
% 利用AI技术自动处理。
% domain-specific task
\noindent
\textbf{Vision applications} on mobile and edge platforms exploit AI technology to process images or videos in RGB spaces~\cite{zhang2018awstream, khani2023recl, Romil2022Ekya}.
Termial devices (\eg smartphones, drones, and wearable cameras) continuously capture massive vision data and transmit to nearby edge servers for analysis~\cite{xu2025joint, dai2024context}.
On the edge server, multiple DNNs that are well-trained on domain-specific datasets separately take care of one target task and together serve the application well, such as image classification~\cite{yuan2022PacketGame, kang2022viva}, vehicle counting~\cite{li2020reducto, kang2017noscope}, and target detection~\cite{du2020server, yu2024vqpy}.
With the natural language interface inherited from LLM, VLMs extend these capabilities to more complex, open-ended vision tasks.
For example, a VLM can locate a target when only given a text-described query such as ``A boy wearing a red sweater lost at the corner''.
In this paradigm, the edge server must handle highly concurrent streams from multiple clients, each with domain-specific tasks.
Instead of deploying multiple full-parameter VLMs, the edge server can utilize LoRA to dynamically equip a shared base model with diverse adapters tailored for different queries.
% multimodal ability and preserve user-friendly

% current video analytics is limited by the poor language understanding capacity of , thus hardly handling multi-modal tasks

% \subsection{Large Multimodal Models}\label{sec:ParaEfficient}
% \vspace{0.2em}
\noindent
\textbf{Vision language models (VLMs)} aim to achieve stronger general intelligence via extending LLMs with multimodal inputs. 
Since more than 80\% of human beings’ perception and activities are mediated through vision~\cite{thomas2024vision}, it is natural to start the exploration by equipping LLMs with ``eyes''.
By introducing a visual receptor, comprised of a \textit{visual encoder} (\eg ViT~\cite{radford2021learning}) and a \textit{vision-language projector} (\eg Q-former~\cite{li2023blip}), VLMs power LLMs with visual capacity. 
Fig. \ref{fig:VLM} illustrates the inference procedure of VLMs.
Given an image input and its prompt, the visual encoder splits the image into small patches (\eg 14$\times$14 pixels block) and extracts the visual features of each. 
The vision-language projector then converts patches' visual features into visual tokens and feeds into LLM~\cite{achiam2023gpt4, chowdhery2023palm, touvron2023llama2, jiang2023mistral}, together with the embedded text tokens from the prompt, to generate the answer. %another sequence of tokens as the answer.
LLM maps the input into high-level features and predicts the probability distribution of the next token with the \textit{language modeling (LM) head} in an autoregressive pattern \cite{kwon2023efficient, yu2022orca}. 
As the dotted arrows depicted in Fig. \ref{fig:VLM}, it iteratively generates tokens with input tokens and previous output tokens once a time, until an end-of-sentence (\texttt{<EOS>}) token is emitted. 

% This parameter-efficient method capitalizes on the observation that only a small portion of model parameters, referred to as adapters, change during fine-tuning~\cite{peft}.


% \subsection{Integrating External Knowledge}\label{sec:ExternalKnowledge}
% 三种方式：few-shot，RAG，LoRA
% 前两种通过修改输入，即增强prompt的方法，后一种通过更改模型参数。
% few-shot怎么做
% RAG怎么做
% 它们的缺点①精度不稳定，且无法保证输出一致；②增加了输入长度，延长推理时延；③RAG的方式需要在大模型外retreival，更耗时。
% 因此，针对有精度和时延要求的vision application中，本文选择LoRA的方式。
% 而对增强prompt的方法的优化放在future work中。


\noindent

\textbf{Low-rank adaptation (LoRA)}~\cite{hu2021lora, dettmers2024qlora, chavan2023one} is a widely used parameter-efficient fine-tuning method~\cite{peft} %, which offers great accuracy gain on 
to integrate the external knowledge into large models.
% by fine-tuning only 
The LoRA adapter is a small number of trainable parameters to learn this knowledge, typically placed in attention layers~\cite{hu2021lora}.
The core of LoRA, as illustrated in Fig. \ref{fig:Alongside}, is to represent each weight update $\Delta W$ as a product of two matrices, $A$ and $B$, with dimensions $r\times d$ and $d\times r$, respectively. 
Their ranks are much smaller than the base model weight $W$, with dimensions $d\times d$.
This makes sense since the low intrinsic rank phenomenon of weight updates in large models~\cite{hu2021lora}.
When fine-tuning, LoRA adapters only update $A$ and $B$ while keeping $W$ frozen. 
At inference, the computation of all LoRA adapters can be placed bypass, as shown in Fig. \ref{fig:Alongside}, and only adds the results onto the output.  
Or merge one multiplied matrix $B\times A$ (\ie $\Delta W$) into base model $W$, which keeps computation overhead the same as the base model in Fig. \ref{fig:Merging}.
% ; such merged inference mode yields the exact inference overhead 


% for instance in Fig. \ref{fig:Alongside}, three adapters 


% \subsection{LoRA Model Serving System}\label{sec:LoRAServing}
% \vspace{0.2em}
\noindent
\textbf{LoRA model serving systems}~\cite{chen2024punica, sheng2023slora, wu2024dlora} aim to improve LoRA model inference efficiency when serving concurrent requests that invokes heterogeneous adapters.
In unmerged inference, the core characteristic is that computidng two small matrices of the adapter underutilizes GPU computing units, leading to unnecessary delay and resource waste.
To tackle this, Punica~\cite{chen2024punica} and S-LoRA~\cite{sheng2023slora} batch multiple heterogeneous adapters, as the cascade LoRA adapters illustrated in Fig. \ref{fig:Alongside}, and compute them within one single custom CUDA kernel.
This method boosts the system throughput indeed but causes significant additional overhead (more in \S \ref{sec:challenges}).
% 增加了吞吐，但是①引入了时延，②没法针对不同时延敏感性的lora分or别计算。
dLoRA~\cite{wu2024dlora} merges the most accessed LoRA adapter into the base model and switches to unmerge if necessary.  % under skewed workloads and unmerge it in other cases.
However, its mode switch causes an unacceptable time cost (more in \S \ref{sec:challenges}).
Besides, like Punica and S-LoRA, dLoRA's unmerged inference also fails to address the large number of extra computation overhead.


